\documentclass[../../main.tex]{subfiles}
\begin{document}

{\bf [解答]}

\paragraph{(1)}

\begin{figure}[htb]
  \centering
  \begin{tikzpicture}[x=1.3cm, y=1.2cm, >=stealth,
    every node/.style={font=\small}]
    \coordinate (O) at (0, 0);
    \coordinate (A) at (1.0, 3.0);
    \coordinate (B) at (5.0, 1.0);
    \coordinate (P) at ($(A)!0.55!(B)$);

    % 直線 AB
    \draw[thick, gray!60] ($(A)!-0.2!(B)$) -- ($(A)!1.2!(B)$)
    node[above right, font=\footnotesize, gray!80!black] {直線 $AB$};

    % 位置ベクトル
    \draw[->, thick, blue!80!black] (O) -- (A) node[midway, above left=2pt] {$\overrightarrow{OA}$};
    \draw[->, thick, blue!80!black] (O) -- (B) node[midway, below right=2pt] {$\overrightarrow{OB}$};
    \draw[->, line width=1.2pt, red!85!black] (O) -- (P) node[midway, above left=2pt] {$\overrightarrow{OP}$};

    % ベクトル AB
    \draw[->, thick, teal!85!black] (A) -- (P) node[midway, above right=1pt, font=\footnotesize] {$t\overrightarrow{AB}$};

    % 各点
    \node[circle, fill, inner sep=1.4pt] at (O) {};
    \node[below left=2pt] at (O) {$O$};
    \node[circle, fill, inner sep=1.4pt] at (A) {};
    \node[above left=2pt] at (A) {$A$};
    \node[circle, fill, inner sep=1.4pt] at (B) {};
    \node[below right=2pt] at (B) {$B$};
    \node[circle, fill=red!85!black, inner sep=1.5pt] at (P) {};
    \node[above right=2pt, red!85!black] at (P) {$P$};
  \end{tikzpicture}
  \caption{直線 $AB$ 上の点 $P$ と位置ベクトル $\overrightarrow{OP}$}
  \label{fig:1967_5_line_ab}
\end{figure}

$P$ が $AB$ 上にある時, 実数 $t$ を用いて
\begin{align}
  \overrightarrow{OP} = \overrightarrow{OA} + t \overrightarrow{AB}
\end{align}
とかける．整理して
\begin{align}
  \overrightarrow{OP} = (1-t)\overrightarrow{OA} + t\overrightarrow{OB}
\end{align}
だから $a = 1-t, b = t$ とおけば, 十分性が示される．

次に, $\overrightarrow{OP} = a\overrightarrow{OA} + b\overrightarrow{OB}, a+b=1$ とおける時, $b=1-a$ から
\begin{align}
  \overrightarrow{OP} &= a\overrightarrow{OA} + (1-a)\overrightarrow{OB} \\
  &= \overrightarrow{OB} + a\overrightarrow{BA}
\end{align}
となり, たしかに $P$ は $AB$ 上にある．よって必要性が示された．

以上から題意は示された．

\paragraph{(2)}

与えられた条件
\begin{align}
  \begin{cases}
    p\overrightarrow{OA}+q\overrightarrow{OB}+r\overrightarrow{OC}=\vec{0} \\
    p+q+r=0
  \end{cases} \label{eq:1}
\end{align}
について考える．第二式から$r=-(p+q)$ だから,
\begin{align}
  & p(\overrightarrow{OA} - \overrightarrow{OC}) + q(\overrightarrow{OB} - \overrightarrow{OC}) = 0 \\
  & p\overrightarrow{CA} + q\overrightarrow{CB} = 0 \label{eq:2}
\end{align}
となる．

ここで, $A, B, C$ は一直線上にないから\cref{eq:2}をみたす $p, q$ は $p=q=0$ のみである．よって\cref{eq:1}から
\begin{align}
  r = 0
\end{align}
となるから $(p, q, r) = (0, 0, 0)$ である．

\paragraph{(3)}

\begin{figure}[htb]
  \centering
  \begin{tikzpicture}[x=1.3cm, y=1.2cm, >=stealth,
    every node/.style={font=\small}]
    \coordinate (O) at (0, 0);
    \coordinate (A) at (3.5, 2.5);
    \coordinate (B) at (6.0, 1.5);
    \coordinate (C) at (2.5, 4.5);

    % AQ = x*AB + y*AC (x=0.8, y=0.2)
    \coordinate (Bx) at ($(A)!0.8!(B)$);
    \coordinate (Cy) at ($(A)!0.2!(C)$);
    \coordinate (Q) at ($(Bx)+(Cy)-(A)$);

    % 平行四辺形補助線
    \draw[dashed, gray!60] (Bx) -- (Q);
    \draw[dashed, gray!60] (Cy) -- (Q);

    % 三角形 ABC
    \draw[thick, gray!50] (A) -- (B) -- (C) -- cycle;

    % 基底ベクトル
    \draw[->, thick, teal!85!black] (A) -- (Bx) node[pos=0.3, below=2pt, font=\footnotesize] {$x\overrightarrow{AB}$};
    \draw[->, thick, teal!85!black] (A) -- (Cy) node[midway, left=2pt, font=\footnotesize] {$y\overrightarrow{AC}$};
    \draw[->, line width=1.1pt, teal!70!black] (A) -- (Q) node[midway, above=2pt, font=\footnotesize] {$\overrightarrow{AQ}$};

    % 位置ベクトル
    \draw[->, thick, blue!80!black] (O) -- (A) node[midway, above left=1pt] {$\overrightarrow{OA}$};
    \draw[->, thick, blue!70!gray] (O) -- (B) node[midway, below right=1pt] {$\overrightarrow{OB}$};
    \draw[->, thick, blue!70!gray] (O) -- (C) node[midway, above left=1pt] {$\overrightarrow{OC}$};
    \draw[->, line width=1.3pt, red!85!black] (O) -- (Q) node[pos=0.45, above left=1pt] {$\overrightarrow{OQ}$};

    % 各点
    \node[circle, fill, inner sep=1.4pt] at (O) {};
    \node[below left=2pt] at (O) {$O$};
    \node[circle, fill, inner sep=1.4pt] at (A) {};
    \node[below=2pt] at (A) {$A$};
    \node[circle, fill, inner sep=1.4pt] at (B) {};
    \node[below right=2pt] at (B) {$B$};
    \node[circle, fill, inner sep=1.4pt] at (C) {};
    \node[above=2pt] at (C) {$C$};
    \node[circle, fill=red!85!black, inner sep=1.5pt] at (Q) {};
    \node[above right=2pt, red!85!black] at (Q) {$Q$};
  \end{tikzpicture}
  \caption{平面 $ABC$ 上の点 $Q$ と位置ベクトル $\overrightarrow{OQ}$}
  \label{fig:1967_5_plane_abc}
\end{figure}

与えられた条件
\begin{align}
  \overrightarrow{AQ}=x\overrightarrow{AB}+y\overrightarrow{AC}
\end{align}
から$\overrightarrow{OQ}$を計算すると
\begin{align}
  \overrightarrow{OQ}
  &= \overrightarrow{OA} + \overrightarrow{AQ} \\
  &= \overrightarrow{OA} + x\overrightarrow{AB} + y\overrightarrow{AC} \\
  &= (1-x-y)\overrightarrow{OA} + x\overrightarrow{OB} + y\overrightarrow{OC}
\end{align}
である．ここで
\begin{align}
  l = 1 - x - y, \quad
  m = x, \quad
  n = y
\end{align}
とおくと，これは$l+m+n=1$を満たしてかつ一意で必ず存在する．
よって題意は示された．

\newpage

\end{document}
